\section{PROCEDIMENTOS METODOLÓGICOS OU MATERIAL E MÉTODOS}
Aqui são descritos os passos metodológicos percorridos pela pesquisa.

\subsection{MODALIDADE DE PESQUISA}
Quanto ao propósito: pesquisa exploratória, descritiva ou explicativa. Quanto aos métodos empregados: segundo a natureza dos dados (quantitativa, qualitativa ou qualiquantitativa).

Informar o ambiente em que os dados foram coletados (pesquisa de campo, laboratório). O tipo de pesquisa: bibliográfica, documental, experimental, ensaio clínico, estudo de coorte, caso-controle, estudo de caso, etnográfica, pesquisa-ação, pesquisa participante, fenomenológica e mista. \cite{gil2017}

A escolha do procedimento metodológico tem relação direta com a necessidade de atingir os objetivos da pesquisa.

\subsection{INSTRUMENTOS DE COLETA DE DADOS}
Informar quais são os instrumentos de coleta de dados (entrevista, aplicação de questionário, levantamento documental etc.). Como se dará a análise dos dados?

Aqui são descritos os passos metodológicos percorridos pela pesquisa.

\subsection{PARTICIPANTES DA PESQUISA}
Deve informar quem são os partícipes da pesquisa, deixando claros os motivos para a escolha daquele público.

\subsubsection{Aspectos éticos da pesquisa}
Avaliar e apresentar as questões éticas que envolvem a realização da pesquisa. Deve informar se há e quais são os riscos na realização da pesquisa para o público participante. Ex.: riscos ergonômicos, psicológicos, físicos etc.